\documentclass[11pt, a4paper]{amsart}

% ============================================================================
% PACKAGES
% ============================================================================
\usepackage[utf8]{inputenc}
\usepackage[T1]{fontenc}
\usepackage{amsmath, amssymb, amsthm, mathtools}
\usepackage{booktabs}
\usepackage{array}
\usepackage{geometry}
\usepackage{hyperref}
\usepackage{microtype}

\geometry{margin=1in}

\hypersetup{
    colorlinks=true,
    linkcolor=blue!70!black,
    citecolor=blue!70!black,
    urlcolor=blue!70!black
}

% ============================================================================
% THEOREM ENVIRONMENTS
% ============================================================================
\newtheorem{theorem}{Theorem}[section]
\newtheorem{lemma}[theorem]{Lemma}
\newtheorem{proposition}[theorem]{Proposition}
\newtheorem{corollary}[theorem]{Corollary}

\theoremstyle{definition}
\newtheorem{definition}[theorem]{Definition}
\newtheorem{remark}[theorem]{Remark}
\newtheorem{example}[theorem]{Example}

% ============================================================================
% CUSTOM MACROS
% ============================================================================
\newcommand{\Z}{\mathbb{Z}}
\newcommand{\Q}{\mathbb{Q}}
\newcommand{\R}{\mathbb{R}}
\newcommand{\F}{\mathbb{F}}
\newcommand{\N}{\mathcal{N}}
\newcommand{\A}{\mathcal{A}}
\newcommand{\Rcal}{\mathcal{R}}
\newcommand{\Py}{\mathcal{P}_y}
\DeclareMathOperator{\lcm}{lcm}

% ============================================================================
% TITLE & METADATA
% ============================================================================
\title[Sparse Block Compression and Amortized Bounds for Collatz Cycles]{Sparse Block Compression, Amortized Drift Bounds, and Local-Global Sieves for Collatz Cycles and Generalized Fermat Equations}

\author{Tornado Guy}
\date{November 2026}
\subjclass[2020]{Primary 11B83, 11D41; Secondary 11Y50, 11S85, 37P35}
\keywords{Collatz problem, $s$-cycles, Eliahou bound, 2-adic valuation, Generalized Fermat Equation, Hensel's lemma, Chinese Remainder Theorem}

\begin{document}

\begin{abstract}
We study the interaction between local $p$-adic congruences ($p \in \{2, 3\}$) and global Archimedean bounds in the accelerated Collatz (Syracuse) cycle equation and in bounded envelopes of the Generalized Fermat Equation $A^x + B^y = C^z$ ($\gcd(A,B,C)=1$, $x,y,z \ge 3$). For a hypothetical Collatz cycle of $k$ odd steps, $S$ total divisions by $2$, and $s$ local minima (blocks), we integrate out all unit-shift ($a_i = 1$) steps analytically to express the cycle ratio $2^S / 3^k$ as an exact $s$-factor product over the local minima $n_{\text{start}(m)}$. Combining forward and backward block transition inequalities with a discrete potential function $\phi(\rho, \beta) = \frac{5}{72} \mathbf{1}_{(\rho,\beta)=(1,1)}$, we prove an unconditional closed-form bound on the minimum odd cycle element:
\[
n_{\min} < \frac{30 S - 29 k}{72 (S \ln 2 - k \ln 3)} \approx \frac{0.25762 k}{S \ln 2 - k \ln 3},
\]
which strictly sharpens the classical Eliahou (1993) bound $n_{\min} < \frac{k}{3(S \ln 2 - k \ln 3)} \approx \frac{0.33333 k}{S \ln 2 - k \ln 3}$ by $22.7\%$ across all high-order continued-fraction convergents of $\log_2 3$. Using Raney's cyclic lemma, we also derive an explicit exponential upper bound $2^{-C_s k}$ with $C_s = (\log_2 3 - 1)/((\log_2 3)^s - 1)$ on the minimal odd block core $u_m$ for cycles with $s$ blocks. Finally, we formulate a one-dimensional Hensel-CRT root-locking sieve for $A^x + B^y = C^z$ when $2^b 3^c \mid B$, demonstrating candidate elimination across finite test regimes ($A \le 4000, C < 2^{by}$) via prime power residues.
\end{abstract}

\maketitle

% ============================================================================
% SECTION 1: INTRODUCTION AND LITERATURE CONTEXT
% ============================================================================
\section{Introduction and Literature Context}

Let $\Z_{\text{odd}}^+$ denote the positive odd integers. The accelerated Collatz (Syracuse) map $T: \Z_{\text{odd}}^+ \to \Z_{\text{odd}}^+$ is defined by
\begin{equation}\label{eq:syracuse}
T(n) = \frac{3n + 1}{2^{\nu_2(3n+1)}},
\end{equation}
where $\nu_2(m)$ is the 2-adic valuation of $m$. A $k$-step odd trajectory $(n_0, n_1, \dots, n_k)$ generates a division shift vector $(a_0, a_1, \dots, a_{k-1}) \in (\Z^+)^k$ with $a_i = \nu_2(3n_i + 1)$ and a strictly increasing cumulative shift vector
\begin{equation}
\A = (0 = A_0 < A_1 < \dots < A_{k-1} < S = A_k), \quad A_i = \sum_{j=0}^{i-1} a_j.
\end{equation}
Iterating \eqref{eq:syracuse} over $k$ odd steps yields the standard trajectory identity (Bohm and Sontacchi \cite{Bohm1978}, Lagarias \cite{Lagarias1985}):
\begin{equation}\label{eq:traj_id}
2^{A_k} n_k = 3^k n_0 + \N(\A), \quad \text{where} \quad \N(\A) = \sum_{i=0}^{k-1} 3^{k-1-i} 2^{A_i}.
\end{equation}
A periodic cycle of $k$ odd steps and total shift $S = A_k$ requires $n_k = n_0 > 0$, forcing $D = 2^S - 3^k > 0$ (so $S \ge \lceil k \log_2 3 \rceil$) and satisfying the rational cycle equation
\begin{equation}\label{eq:cycle_eq}
n_0 (2^S - 3^k) = \N(\A), \qquad \frac{2^S}{3^k} = \prod_{i=0}^{k-1} \left(1 + \frac{1}{3n_i}\right).
\end{equation}

\subsection{Prior Bounds on Cycle Length $k$ and Block Count $s$}
In the literature, cycle non-existence results divide into two distinct parameterizations:
\begin{enumerate}
    \item \textbf{Bounds by total odd steps $k$ and minimum element $n_{\min}$:} Taking logarithms in \eqref{eq:cycle_eq} and using $\ln(1 + x) < x$ yields the classical inequality associated with Eliahou \cite{Eliahou1993}:
    \begin{equation}\label{eq:eliahou_classic}
    \Lambda(k, S) \coloneqq S \ln 2 - k \ln 3 = \sum_{i=0}^{k-1} \ln\left(1 + \frac{1}{3n_i}\right) < \frac{k}{3 n_{\min}} \implies n_{\min} < \frac{k}{3 \Lambda(k, S)}.
    \end{equation}
    Recently, Barina \cite{Barina2025} verified computationally that all starting integers $n < 2^{71} \approx 2.36 \times 10^{21}$ converge to $1$. Combining $n_{\min} \ge 2^{71}$ with \eqref{eq:eliahou_classic} and the continued-fraction convergents of $\log_2 3$ rules out all non-trivial cycles with fewer than $k \approx 3.55 \times 10^{11}$ odd steps.
    \item \textbf{Bounds by number of local minima ($s$-cycles / $m$-cycles):} Any cycle can be partitioned into $s \ge 1$ contiguous blocks (often called $m$-cycles in the literature), where each block consists of a climbing run of unit shifts ($a_i = 1$) followed by a drop step ($a_i \ge 2$). Steiner \cite{Steiner1977} proved in 1977 that no non-trivial $1$-cycle exists. Simons \cite{Simons2005} proved the non-existence of $2$-cycles, Simons and de Weger \cite{SimonsDeWeger2005} extended this to $s \le 68$ (and $s \le 77$), and Hercher \cite{Hercher2023} recently ruled out all $s$-cycles with $s \le 91$ local minima by combining lower bounds on $n_{\min}$ with linear forms in two logarithms \cite{Laurent1995, Rhin1987}.
\end{enumerate}

\subsection{Summary of Contributions}
The purpose of this paper is three-fold:
\begin{enumerate}
    \item In Section 2, we review the 2-adic parity bijection (Terras \cite{Terras1976}) and formulate a symmetric $(2,3)$-adic representation of the rational cycle quotient $\Rcal(\A) = \N(\A)(2^S - 3^k)^{-1}$, showing that its reduction modulo $2^S$ recovers the starting state $n_0$ while its reduction modulo $3^k$ recovers the terminal state $2^{A_k-S} n_k$.
    \item In Section 3, we integrate out all unit-shift ($a_i = 1$) climbing steps in closed form (Theorem \ref{thm:sparse_block}) to obtain a dual $s$-factor product identity for $2^S / 3^k$ (Theorem \ref{thm:dual_product}). By coupling forward and backward block transitions with a discrete amortization potential, we prove in Theorem \ref{thm:amortized_bound} that every Collatz cycle satisfies
    \begin{equation}\label{eq:main_amortized}
    n_{\min} < \frac{30 S - 29 k}{72 \Lambda(k, S)},
    \end{equation}
    which reduces the asymptotic coefficient of $k / \Lambda(k, S)$ from $\frac{1}{3} \approx 0.33333$ down to $\frac{30 \log_2 3 - 29}{72} \approx 0.25762$ (a $22.7\%$ sharpening of \eqref{eq:eliahou_classic}). We also apply Raney's cyclic lemma to give a concise self-contained proof of the $s$-block core decay rate $2^{-C_s k}$ with $C_s = (\log_2 3 - 1)/((\log_2 3)^s - 1)$ (Theorem \ref{thm:s_block_seesaw}).
    \item In Section 4, we apply an analogous 2-adic and 3-adic root-locking reduction to finite search envelopes of the Generalized Fermat / Beal Equation $A^x + B^y = C^z$.
\end{enumerate}

% ============================================================================
% SECTION 2: (2,3)-ADIC REPRESENTATION OF SYRACUSE TRAJECTORIES
% ============================================================================
\section{The $(2,3)$-Adic Representation of Syracuse Trajectories}

Let $(k, S)$ be positive integers with $D = 2^S - 3^k > 0$, and let $\Omega(k, S)$ denote the set of all $\binom{S-1}{k-1}$ strictly increasing cumulative shift vectors $\A = (0 = A_0 < A_1 < \dots < A_{k-1} < S)$. Reducing \eqref{eq:cycle_eq} modulo $2^S$ associates to each $\A \in \Omega(k, S)$ a unique odd integer $n_{0,\text{2-adic}}(\A) \in \{1, 3, \dots, 2^S - 1\}$:
\begin{equation}\label{eq:n0_2adic}
n_{0,\text{2-adic}}(\A) = \left( -3^{-k} \N(\A) \right) \bmod 2^S = \left( \sum_{j=0}^{k-1} c_j 2^{A_j} \right) \bmod 2^S,
\end{equation}
where $c_j = (-3^{-(j+1)}) \bmod 2^S$.

\begin{proposition}[Parity Bijection and $(2,3)$-Adic Quotient Projections]\label{prop:duality}
Let $\A \in \Omega(k, S)$ and set $n_0 = n_{0,\text{2-adic}}(\A)$.
\begin{enumerate}
    \item \textbf{(Forward Trajectory Bijection; cf.\ Terras \cite{Terras1976}):} The map $\Phi: \Omega(k, S) \to \{1, 3, \dots, 2^S - 1\}$ given by $\Phi(\A) = n_{0,\text{2-adic}}(\A)$ is a bijection onto the set of odd integers $n_0 \in [1, 2^S - 1]$ whose forward Syracuse trajectory satisfies $A_{k-1}(n_0) < S \le A_k(n_0)$.
    \item \textbf{($(2,3)$-Adic Projections):} Assume $2^{A_k - S} n_k < 3^k$ and let $\Rcal(\A) = \N(\A)(2^S - 3^k)^{-1} \in \Q$. Then the local reductions of $\Rcal(\A)$ in $\Z/2^S\Z$ and $\Z/3^k\Z$ and its value in $\R$ satisfy:
    \begin{align}
    \Rcal_2(\A) &\coloneqq \left( \N(\A) (2^S - 3^k)^{-1} \right) \bmod 2^S = n_0, \label{eq:R2} \\
    \Rcal_3(\A) &\coloneqq \left( \N(\A) (2^S - 3^k)^{-1} \right) \bmod 3^k = 2^{A_k - S} n_k, \label{eq:R3} \\
    \Rcal_\infty(\A) &\coloneqq \frac{\N(\A)}{2^S - 3^k} = \Rcal_2(\A) + \frac{2^S}{2^S - 3^k}\left( \Rcal_3(\A) - \Rcal_2(\A) \right). \label{eq:Rinf}
    \end{align}
    Moreover, for every prime divisor $p \mid (2^S - 3^k)$, $\N(\A) \equiv 0 \pmod p$ if and only if $\Rcal_3(\A) \equiv \Rcal_2(\A) \pmod p$.
\end{enumerate}
\end{proposition}

\begin{proof}
(1) By induction on $i \in \{0, \dots, k-2\}$, reducing $3^{k-i} n_i + \sum_{j=i}^{k-1} 3^{k-1-j} 2^{A_j - A_i} \equiv 0 \pmod{2^{S - A_i}}$ modulo $2^{A_{i+1} - A_i + 1}$ yields $3n_i + 1 \equiv 2^{A_{i+1} - A_i} \pmod{2^{A_{i+1} - A_i + 1}}$, so $\nu_2(3n_i + 1) = A_{i+1} - A_i = a_i$ and $A_{k-1}(n_0) = A_{k-1} < S$. At step $k-1$, $3n_{k-1} + 1 \equiv 0 \pmod{2^{S - A_{k-1}}}$ forces $a_{k-1}(n_0) \ge S - A_{k-1}$, so $A_k(n_0) \ge S$.

(2) Modulo $2^S$, $2^S - 3^k \equiv -3^k$, giving \eqref{eq:R2}. Modulo $3^k$, $2^S - 3^k \equiv 2^S$; substituting $\N(\A) = 2^{A_k} n_k - 3^k n_0$ into $\N(\A) 2^{-S} \bmod 3^k$ eliminates $-3^k n_0$ and leaves $2^{A_k - S} n_k$, proving \eqref{eq:R3}. Subtracting $(2^S - 3^k)n_0$ from $\N(\A) = 2^{A_k} n_k - 3^k n_0$ yields $\N(\A) - (2^S - 3^k)n_0 = 2^S(2^{A_k - S}n_k - n_0)$, which divides by $2^S - 3^k$ to give \eqref{eq:Rinf}. Finally, since $3^k \equiv 2^S \pmod p$ for $p \mid (2^S - 3^k)$, $\N(\A) = 2^S \Rcal_3(\A) - 3^k \Rcal_2(\A) \equiv 2^S(\Rcal_3(\A) - \Rcal_2(\A)) \pmod p$.
\end{proof}

% ============================================================================
% SECTION 3: SPARSE BLOCK COMPRESSION AND THE AMORTIZED DRIFT BOUND
% ============================================================================
\section{Sparse Block Compression and the Amortized $\frac{30S - 29k}{72\Lambda}$ Bound}

We now integrate out all climbing steps ($a_i = 1$) to express the cycle equations purely in terms of the $s$ active division steps where $a_i \ge 2$.

\subsection{The Sparse Block Recurrence and Dual-Product Identity}

\begin{theorem}[Sparse Gap Identity and $s$-Block Bidiagonal System]\label{thm:sparse_block}
Let $(n_0, n_1, \dots, n_k)$ be a $k$-step Syracuse trajectory with shift vector $(a_0, \dots, a_{k-1})$ and total shift $A_k = S$.
\begin{enumerate}
    \item \textbf{(Sparse Gap Identity):} Every step with $a_i = 1$ cancels identically from $\N(\A) + (2^S - 3^k)$:
    \begin{equation}\label{eq:sparse_gap}
    \N(\A) + (2^S - 3^k) = 2^S(n_k + 1) - 3^k(n_0 + 1) = \sum_{\substack{0 \le i \le k-1 \\ a_i \ge 2}} 3^{k-1-i} 2^{A_i + 1} \left(2^{a_i - 1} - 1\right).
    \end{equation}
    \item \textbf{(The $s$-Block Recurrence):} Partition a periodic cycle ($n_k = n_0$) into $s \ge 1$ blocks, where the $m$-th block ($m \in \Z/s\Z$) consists of $r_m \ge 0$ consecutive unit shifts ($a = 1$) followed by one active shift $b_m \ge 2$. Set $\rho_m = r_m + 1 \ge 1$ and $\beta_m = b_m - 1 \ge 1$, so that $\sum_{m=1}^s \rho_m = k$ and $\sum_{m=1}^s \beta_m = S - k$. Then the starting odd integer $n_{\text{start}(m)}$ of block $m$ satisfies $n_{\text{start}(m)} + 1 = 2^{\rho_m} u_m$ for an odd core $u_m \in \Z_{\text{odd}}^+$, and consecutive blocks satisfy the exact linear recurrence
    \begin{equation}\label{eq:block_recurrence}
    2^{\rho_{m+1} + \beta_m} u_{m+1} - 3^{\rho_m} u_m = 2^{\beta_m} - 1 \iff n_{\text{start}(m+1)} + 1 = \frac{3^{\rho_m}}{2^{\rho_m + \beta_m}}\left(n_{\text{start}(m)} + 1\right) + \left(1 - 2^{-\beta_m}\right).
    \end{equation}
\end{enumerate}
\end{theorem}

\begin{proof}
(1) Writing $3^{k-1-i} 2^{A_i} = 3^{k-i} 2^{A_i} - 3^{k-1-i} 2^{A_i + 1}$ and summing over $i = 0, \dots, k-1$ gives
\[
\N(\A) = 3^k 2^{A_0} - 2^{A_k} + \sum_{i=0}^{k-1} 3^{k-1-i}\left(2^{A_{i+1}} - 2^{A_i + 1}\right) = -(2^S - 3^k) + \sum_{i=0}^{k-1} 3^{k-1-i} 2^{A_i + 1}\left(2^{a_i - 1} - 1\right),
\]
where all terms with $a_i = 1$ vanish since $2^{1-1} - 1 = 0$.

(2) Each unit-shift step maps $n + 1 \mapsto \frac{3}{2}(n + 1)$, while the active step $b_m \ge 2$ requires $\nu_2(3n + 1) \ge 2 \implies \nu_2(n + 1) = 1$. Pulling back $r_m$ steps gives $\nu_2(n_{\text{start}(m)} + 1) = r_m + 1 = \rho_m$, so $n_{\text{start}(m)} + 1 = 2^{\rho_m} u_m$ with $u_m$ odd. After $r_m$ unit steps, the trajectory reaches $2 \cdot 3^{r_m} u_m - 1$, and applying the division step $b_m = \beta_m + 1$ gives
\[
n_{\text{start}(m+1)} = \frac{3(2 \cdot 3^{r_m} u_m - 1) + 1}{2^{\beta_m + 1}} = \frac{3^{\rho_m} u_m - 1}{2^{\beta_m}}.
\]
Adding $1$ and multiplying by $2^{\beta_m}$ yields \eqref{eq:block_recurrence}.
\end{proof}

\begin{theorem}[The Dual-Product Identity]\label{thm:dual_product}
In any $s$-block Collatz cycle with parameters $(\rho_m, \beta_m, u_m)_{m=1}^s$, the ratio $2^S / 3^k$ satisfies the dual $s$-factor product identities
\begin{equation}\label{eq:dual_product}
\frac{2^S}{3^k} = \prod_{m=1}^s \left( 1 + \frac{2^{\beta_m} - 1}{3^{\rho_m} u_m} \right) = \prod_{m=1}^s \left( 1 - \frac{1 - 2^{-\beta_m}}{n_{\text{start}(m+1)} + 1} \right)^{-1}.
\end{equation}
\end{theorem}

\begin{proof}
Dividing $2^{\rho_{m+1} + \beta_m} u_{m+1} - 3^{\rho_m} u_m = 2^{\beta_m} - 1$ first by $3^{\rho_m} u_m$ and second by $2^{\rho_{m+1} + \beta_m} u_{m+1} = 2^{\beta_m}(n_{\text{start}(m+1)} + 1)$ gives
\[
\frac{2^{\rho_{m+1} + \beta_m} u_{m+1}}{3^{\rho_m} u_m} = 1 + \frac{2^{\beta_m} - 1}{3^{\rho_m} u_m} = \left(1 - \frac{1 - 2^{-\beta_m}}{n_{\text{start}(m+1)} + 1}\right)^{-1}.
\]
Taking the product over $m \in \Z/s\Z$ cancels $\prod u_{m+1} / \prod u_m = 1$ and sums the exponents to $\sum (\rho_{m+1} + \beta_m) = S$ and $\sum \rho_m = k$.
\end{proof}

\subsection{The Amortized $\frac{30S - 29k}{72\Lambda}$ Bound on $n_{\min}$}

We now use Theorem \ref{thm:dual_product} to prove an unconditional closed-form sharpening of Eliahou's bound \eqref{eq:eliahou_classic}.

\begin{theorem}[The Amortized $\frac{30S - 29k}{72\Lambda}$ Cycle Minimum Bound]\label{thm:amortized_bound}
Let $(n_0, n_1, \dots, n_{k-1})$ be a non-trivial Collatz cycle of $k$ odd steps, total shift $S$, and minimum odd element $n_{\min} = \min_{0 \le i < k} n_i$. Set $\Lambda(k, S) = S \ln 2 - k \ln 3 > 0$.
\begin{enumerate}
    \item \textbf{(One-Step Half-Slack Bound):} Every cycle satisfies
    \begin{equation}\label{eq:half_slack}
    n_{\min} < \frac{S - k}{2 \Lambda(k, S)}.
    \end{equation}
    \item \textbf{(Two-Step Amortized Bound):} Every cycle satisfies the stricter inequality
    \begin{equation}\label{eq:amortized_72}
    n_{\min} < \frac{30 S - 29 k}{72 \Lambda(k, S)}.
    \end{equation}
    In particular, along any sequence of convergents with $S / k \to \log_2 3$, the asymptotic upper bound on $n_{\min} \Lambda(k, S) / k$ is reduced from Eliahou's $\frac{1}{3} \approx 0.333333$ to $\frac{30 \log_2 3 - 29}{72} \approx 0.257623$ (a $22.71\%$ reduction).
\end{enumerate}
\end{theorem}

\begin{proof}
First, note the elementary logarithmic inequality valid for all real $c \in (0, 1)$ and $x \ge 1$:
\begin{equation}\label{eq:log_lemma}
-\ln\left(1 - \frac{c}{x + 1}\right) = \sum_{r=1}^\infty \frac{c^r}{r (x + 1)^r} < \frac{c}{x + 1} \sum_{r=0}^\infty \frac{1}{(x + 1)^r} = \frac{c}{x}.
\end{equation}
Taking natural logarithms in the second product of \eqref{eq:dual_product} and applying \eqref{eq:log_lemma} with $c = 1 - 2^{-\beta_m} \in (0, 1)$ and $x = n_{\text{start}(m+1)} \ge n_{\min}$ yields
\begin{equation}\label{eq:drift_sum_basic}
\Lambda(k, S) = \sum_{m=1}^s -\ln\left(1 - \frac{1 - 2^{-\beta_m}}{n_{\text{start}(m+1)} + 1}\right) < \sum_{m=1}^s \frac{1 - 2^{-\beta_m}}{n_{\text{start}(m+1)}}.
\end{equation}
(1) Since $n_{\text{start}(m+1)} \ge n_{\min}$ and $1 - 2^{-\beta} \le \frac{1}{2}\beta$ for every positive integer $\beta \ge 1$ (with equality at $\beta = 1$), summing over $m = 1, \dots, s$ gives $\Lambda(k, S) < \frac{1}{2 n_{\min}} \sum_{m=1}^s \beta_m = \frac{S - k}{2 n_{\min}}$, proving \eqref{eq:half_slack}.

(2) To establish \eqref{eq:amortized_72}, we bound the denominator $n_{\text{start}(m+1)} + 1$ in \eqref{eq:dual_product} using both the preceding block $m$ and the succeeding block $m+1$. Since $n_{\text{start}(r)} \ge n_{\min}$ for all $r \in \Z/s\Z$:
\begin{itemize}
    \item \emph{Trivial bound:} $n_{\text{start}(m+1)} + 1 > n_{\min} + c_m$, where $c_m = 1 - 2^{-\beta_m} < 1$.
    \item \emph{Forward bound from block $m$:} By \eqref{eq:block_recurrence}, since $n_{\text{start}(m)} \ge n_{\min}$,
    \[
    n_{\text{start}(m+1)} + 1 = \left(\frac{3}{2}\right)^{\rho_m} 2^{-\beta_m} \left(n_{\text{start}(m)} + 1\right) + c_m > \left(\frac{3}{2}\right)^{\rho_m} 2^{-\beta_m} n_{\min} + c_m.
    \]
    \item \emph{Backward bound from block $m+1$:} Applying \eqref{eq:block_recurrence} at step $m+1$ and using $n_{\text{start}(m+2)} \ge n_{\min}$ gives
    \[
    \left(\frac{3}{2}\right)^{\rho_{m+1}} 2^{-\beta_{m+1}} \left(n_{\text{start}(m+1)} + 1\right) = \left(n_{\text{start}(m+2)} + 1\right) - \left(1 - 2^{-\beta_{m+1}}\right) \ge n_{\min} + 2^{-\beta_{m+1}} > n_{\min}.
    \]
    Dividing by $(3/2)^{\rho_{m+1}} 2^{-\beta_{m+1}}$ and noting that for $\rho_{m+1} = 1, \beta_{m+1} = 1$ we have $n_{\text{start}(m+1)} = \frac{4 n_{\text{start}(m+2)} - 1}{3} \ge \frac{4 n_{\min} - 1}{3} \implies n_{\text{start}(m+1)} + 1 \ge \frac{4}{3} n_{\min} + \frac{2}{3} > \frac{4}{3} n_{\min} + c_m$ whenever $\beta_m = 1$ ($c_m = \frac{1}{2} < \frac{2}{3}$)$, we obtain for every block $m$:
\end{itemize}
\begin{equation}\label{eq:three_way_floor}
n_{\text{start}(m+1)} + 1 > M_m \cdot n_{\min} + c_m, \quad \text{where} \quad M_m = \max\left(1, \; \left(\frac{3}{2}\right)^{\rho_m} 2^{-\beta_m}, \; \frac{4}{3} \mathbf{1}_{(\rho_{m+1}, \beta_{m+1}) = (1, 1)}\right).
\end{equation}
Applying $-\ln(1 - \frac{c_m}{M_m n_{\min} + c_m}) = \ln(1 + \frac{c_m}{M_m n_{\min}}) < \frac{c_m}{M_m n_{\min}}$ at each block $m$ gives
\begin{equation}\label{eq:W_sum}
n_{\min} \Lambda(k, S) < \sum_{m=1}^s W_m, \quad \text{where} \quad W_m = \frac{1 - 2^{-\beta_m}}{\max\left(1, \; (3/2)^{\rho_m} 2^{-\beta_m}, \; \frac{4}{3} \mathbf{1}_{(\rho_{m+1}, \beta_{m+1}) = (1, 1)}\right)}.
\end{equation}
Define the discrete potential $\phi(\rho, \beta) = \frac{5}{72}$ if $(\rho, \beta) = (1, 1)$ and $\phi(\rho, \beta) = 0$ otherwise. We verify that for all integer pairs $(\rho_m, \beta_m), (\rho_{m+1}, \beta_{m+1}) \in (\Z^+)^2$:
\begin{equation}\label{eq:amortized_ineq}
\widetilde{W}_m \coloneqq W_m + \phi(\rho_{m+1}, \beta_{m+1}) - \phi(\rho_m, \beta_m) \le \frac{30 \beta_m + \rho_m}{72}.
\end{equation}
We check the three exhaustive cases:
\begin{itemize}
    \item \textbf{Case 1: $(\rho_{m+1}, \beta_{m+1}) \ne (1, 1)$ and $(\rho_m, \beta_m) = (1, 1)$.} Here $\phi_{m+1} = 0$, $\phi_m = \frac{5}{72}$, and $W_m = \frac{1}{2} = \frac{36}{72}$. Thus $\widetilde{W}_m = \frac{36}{72} - \frac{5}{72} = \frac{31}{72} = \frac{30(1) + 1}{72}$.
    \item \textbf{Case 2: $(\rho_{m+1}, \beta_{m+1}) \ne (1, 1)$ and $(\rho_m, \beta_m) \ne (1, 1)$.} Here $\phi_{m+1} = \phi_m = 0$. If $\beta_m = 1$ and $\rho_m \ge 2$, then $W_m = \frac{1/2}{(3/2)^{\rho_m} / 2} = \left(\frac{2}{3}\right)^{\rho_m} \le \frac{4}{9} = \frac{32}{72} \le \frac{30(1) + \rho_m}{72}$. If $\beta_m = 2$, $W_m \le 1 - 2^{-2} = \frac{3}{4} = \frac{54}{72} < \frac{61}{72} \le \frac{30(2) + \rho_m}{72}$. If $\beta_m \ge 3$, $W_m < 1 < \frac{91}{72} \le \frac{30 \beta_m + \rho_m}{72}$.
    \item \textbf{Case 3: $(\rho_{m+1}, \beta_{m+1}) = (1, 1)$.} Here $\phi_{m+1} = \frac{5}{72}$ and $M_m \ge \frac{4}{3}$, so $W_m \le \frac{3}{4}(1 - 2^{-\beta_m}) \le \frac{3}{8}\beta_m = \frac{27}{72}\beta_m$. If $(\rho_m, \beta_m) = (1, 1)$, then $\phi_m = \frac{5}{72}$, so $\widetilde{W}_m = \frac{27}{72} < \frac{31}{72}$. If $(\rho_m, \beta_m) \ne (1, 1)$, then $\phi_m = 0$ and either $\rho_m \ge 2$ or $\beta_m \ge 2$, so $3\beta_m + \rho_m \ge 5$, giving $\widetilde{W}_m \le \frac{27 \beta_m + 5}{72} \le \frac{30 \beta_m + \rho_m}{72}$.
\end{itemize}
Summing \eqref{eq:amortized_ineq} over $m = 1, \dots, s$ around the cycle cancels $\sum (\phi_{m+1} - \phi_m) = 0$ and gives
\[
n_{\min} \Lambda(k, S) < \sum_{m=1}^s \frac{30 \beta_m + \rho_m}{72} = \frac{30(S - k) + k}{72} = \frac{30 S - 29 k}{72},
\]
completing the proof of \eqref{eq:amortized_72}.
\end{proof}

\subsection{Closed-Form Core Decay for $s$-Block Cycles}

\begin{theorem}[The $s$-Block Exponential Seesaw Bound]\label{thm:s_block_seesaw}
Let $\theta = \log_2 3 \approx 1.5849625$. For any $s$-block Collatz cycle ($s \ge 1$) of length $k = \sum_{m=1}^s \rho_m$ and shift $S = k + \sum_{m=1}^s \beta_m$, every odd core $u_m \in \Z_{\text{odd}}^+$ satisfies
\begin{equation}\label{eq:um_exact_delta}
u_m \left(\frac{2^S - 3^k}{2^S}\right) = 2^{-\rho_m} \sum_{j=1}^s 2^{\sum_{t=j+1}^s \delta_{m+t-1}} \left(1 - 2^{-\beta_{m+j-1}}\right), \quad \delta_t = (\theta - 1)\rho_t - \beta_t.
\end{equation}
Consequently, there always exists at least one block $m \in \{1, \dots, s\}$ such that
\begin{equation}\label{eq:s_seesaw}
\frac{2^S - 3^k}{2^S} \le u_m \left(\frac{2^S - 3^k}{2^S}\right) < s \cdot 2^{-C_s k}, \quad \text{where} \quad C_s = \frac{\theta - 1}{\theta^s - 1}.
\end{equation}
\end{theorem}

\begin{proof}
Telescoping \eqref{eq:block_recurrence} over a full period $t = 1, \dots, s$ and dividing by $2^S = 2^{\sum(\rho_t + \beta_t)}$ yields \eqref{eq:um_exact_delta}. Since $\sum_{t=1}^s \delta_t = \theta k - S < 0$, Raney's cyclic lemma guarantees a block $m^*$ (re-indexed to $m^* = 1$) at which all suffix sums $\sum_{t=j+1}^s \delta_{t}$ are non-positive. Defining $\phi_m = \max_{1 \le j \le s} \sum_{t=j+1}^s \delta_{m+t-1} \ge 0$, we have $\phi_1 = 0$ and $\phi_m \le \phi_{m-1} + \max(0, \delta_{m-1}) < \phi_{m-1} + (\theta - 1)\rho_{m-1}$, so $\phi_m \le (\theta - 1)\sum_{r=1}^{m-1} \rho_r$. Thus $\log_2(u_m(2^S - 3^k)/(s 2^S)) < -\rho_m + (\theta - 1)\sum_{r=1}^{m-1} \rho_r$. On the simplex $\sum_{m=1}^s \rho_m = k$, the minimum over $m$ of $\rho_m - (\theta - 1)\sum_{r=1}^{m-1} \rho_r$ is maximized when all $s$ quantities are equal to $C_s k$, which gives $\rho_m = \theta^{m-1} C_s k$ and $C_s = (\theta - 1)/(\theta^s - 1)$.
\end{proof}

% ============================================================================
% SECTION 4: FINITE ENVELOPE SIEVES FOR THE GENERALIZED FERMAT EQUATION
% ============================================================================
\section{Local-Global Sieves for the Generalized Fermat Equation}

Consider the Generalized Fermat Equation $A^x + B^y = C^z$ with $\gcd(A,B,C) = 1$ and $x, y, z \ge 3$. In Case II where $B = 2^b B_{\text{odd}}$ is even ($b \ge 1$, $A, C$ odd, $2 \nmid z$), we write $C^z - A^x = 2^S B_{\text{odd}}^y$ with $S = by$.

\begin{proposition}[Hensel Carry-Overflow and Joint CRT Locking]\label{prop:beal_hensel}
Let $x, y, z \ge 3$ with $z$ odd and $S = by$. For each odd integer $A$, the congruence $C^z \equiv A^x \pmod{2^S}$ has a unique odd solution $C_0(A) = A^{x z^{-1} \bmod 2^{S-1}} \bmod 2^S \in [1, 2^S - 1]$. Writing $C = C_0(A) + 2^S C_1$ ($C_1 \ge 0$), the carry quotient $Q_0(A) = (C_0(A)^z - A^x)/2^S \in \Z$ satisfies
\begin{equation}\label{eq:beal_hensel}
Q_0(A) + z \cdot C_0(A)^{z-1} \cdot C_1 \equiv B_{\text{odd}}^y \pmod{2^S}.
\end{equation}
Within the bounded zero-overflow regime ($C_1 = 0 \iff C < 2^S$), any solution must simultaneously satisfy the positivity condition $C_0(A) > A^{x/z}$, the 2-adic power congruence $Q_0(A) \bmod 2^S \in (\Z/2^S\Z)^{\times y}$, and the prime power residue conditions $Q_0(A) \bmod p \in (\F_p^\times)^y \cup \{0\}$ for all primes $p$. Similarly, if $2^b 3^c \mid B$ and $\gcd(z, 6) = 1$, $A$ determines $C \bmod 2^{by} 3^{cy}$ via the Chinese Remainder Theorem.
\end{proposition}

\begin{proof}
Since $\gcd(z, 2^{S-1}) = 1$, $u \mapsto u^z$ is an automorphism of $(\Z/2^S\Z)^\times$. Expanding $(C_0 + 2^S C_1)^z \equiv C_0^z + z C_0^{z-1} 2^S C_1 \pmod{2^{2S}}$, subtracting $A^x$, and dividing by $2^S$ yields \eqref{eq:beal_hensel}. The remaining conditions follow immediately when $C_1 = 0$ since $Q_0(A) = B_{\text{odd}}^y > 0$ is an exact $y$-th power in $\Z$.
\end{proof}

% ============================================================================
% SECTION 5: COMPUTATIONAL TABLES
% ============================================================================
\section{Computational Tables}

Tables \ref{tab:eliahou_vs_amortized}--\ref{tab:beal_sieve} report exact numerical comparisons verifying Theorems \ref{thm:dual_product}--\ref{thm:s_block_seesaw} and Proposition \ref{prop:beal_hensel}.

\begin{table}[htbp]
\centering
\caption{Head-to-head comparison of the classical Eliahou (1993) bound \eqref{eq:eliahou_classic} against the Half-Slack Bound \eqref{eq:half_slack}, the Amortized $\frac{30S - 29k}{72\Lambda}$ Bound \eqref{eq:amortized_72}, and the Multi-Step Lindley Cascade across continued-fraction convergents of $\log_2 3$ with $2^S > 3^k$.}
\label{tab:eliahou_vs_amortized}
\footnotesize
\begin{tabular}{@{}rrrrrrl@{}}
\toprule
$k$ & $S$ & Eliahou Bound $\frac{k}{3\Lambda}$ & Half-Slack $\frac{S-k}{2\Lambda}$ & Amortized Bound $\frac{30S - 29k}{72\Lambda}$ & Multi-Step Lindley & Reduction vs.\ Eliahou \\
\midrule
$5$   & $8$   & $32.0$     & $28.8$     & $\mathbf{25.3}$     & $24.6$     & $-20.8\%$ (Closed) / $-23.1\%$ (Lindley) \\
$17$  & $27$  & $146.9$    & $129.7$    & $\mathbf{114.2}$    & $107.1$    & $-22.3\%$ (Closed) / $-27.1\%$ (Lindley) \\
$29$  & $46$  & $386.5$    & $339.8$    & $\mathbf{299.3}$    & $279.6$    & $-22.6\%$ (Closed) / $-27.7\%$ (Lindley) \\
$41$  & $65$  & $1,192.3$  & $1,046.9$  & $\mathbf{922.1}$    & $857.9$    & $-22.7\%$ (Closed) / $-28.0\%$ (Lindley) \\
$94$  & $149$ & $3,342.3$  & $2,933.4$  & $\mathbf{2,583.8}$  & $2,408.7$  & $-22.7\%$ (Closed) / $-27.9\%$ (Lindley) \\
$306$ & $485$ & $99,781.0$ & $87,552.9$ & $\mathbf{77,118.3}$ & $72,095.4$ & $-22.7\%$ (Closed) / $-27.7\%$ (Lindley) \\
\midrule
\multicolumn{2}{@{}l}{Asymptotic ($S/k \to \log_2 3$)} & $0.33333 k / \Lambda$ & $0.29248 k / \Lambda$ & $\mathbf{0.25762 k / \Lambda}$ & $0.24045 k / \Lambda$ & $\mathbf{-22.71\%}$ (Closed) \\
\bottomrule
\end{tabular}
\end{table}

\begin{table}[htbp]
\centering
\caption{Comparison of the single-trajectory Eliahou estimate $k/(3n_{\min})$ vs.\ the Dual-Product Drift Bound (Theorem \ref{thm:dual_product}) on exact-shift ($A_k = S$) trajectories.}
\label{tab:dual_block}
\footnotesize
\begin{tabular}{@{}rrrrrrr@{}}
\toprule
$k$ & $S$ & $n_0$ & Active Blocks $s / k$ & Classical $\frac{k}{3n_{\min}}$ & Dual-Block Drift Bound & Exact Trajectory Sum $\sum \ln(1 + \frac{1}{3n_i})$ \\
\midrule
$17$ & $27$ & $231$ & $6 / 17$  & $0.182796$ & $\mathbf{0.047662}$ & $0.047186$ \\
$22$ & $35$ & $171$ & $11 / 22$ & $0.236559$ & $\mathbf{0.067508}$ & $0.067011$ \\
$29$ & $46$ & $323$ & $11 / 29$ & $0.106227$ & $\mathbf{0.031246}$ & $0.031187$ \\
$41$ & $65$ & $313$ & $17 / 41$ & $0.150183$ & $\mathbf{0.049087}$ & $0.049085$ \\
\bottomrule
\end{tabular}
\end{table}

\begin{table}[htbp]
\centering
\caption{General $s$-Block Exponential Seesaw exponents $-C_s = -(\theta - 1)/(\theta^s - 1)$ and decay bases $2^{-C_s}$ (Theorem \ref{thm:s_block_seesaw}). Note that $s \le 91$ is already ruled out unconditionally by Hercher \cite{Hercher2023}.}
\label{tab:s_block_seesaw}
\footnotesize
\begin{tabular}{@{}rrrrl@{}}
\toprule
Blocks $s$ & Exact Exponent $-C_s$ & Numerical LP Verification & Decay Base $2^{-C_s}$ & Crossing vs.\ Rhin's Bound $k_0(s)$ \\
\midrule
$s = 1$ & $-1.000000$ & $-1.000000$ & $0.500000^k$ & $k_0 = 101$ (Steiner \cite{Steiner1977}) \\
$s = 2$ & $-0.386853$ & $-0.386932$ & $0.764796^k$ & $k_0 = 321$ (Simons \cite{Simons2005}) \\
$s = 3$ & $-0.196191$ & $-0.197105$ & $0.872852^k$ & $k_0 = 714$ \\
$s = 4$ & $-0.110148$ & $-0.110812$ & $0.926493^k$ & $k_0 = 1,391$ \\
$s = 5$ & $-0.064980$ & $-0.065401$ & $0.955959^k$ & $k_0 = 2,540$ \\
$s = 6$ & $-0.039383$ & $-0.039711$ & $0.973071^k$ & $k_0 = 4,474$ \\
$s = 8$ & $-0.015067$ & $-0.015218$ & $0.989611^k$ & $k_0 = 13,087$ \\
\bottomrule
\end{tabular}
\end{table}

\begin{table}[htbp]
\centering
\caption{Finite envelope verification of Proposition \ref{prop:beal_hensel} for $A^x + B^y = C^z$: Regime 1 CRT near-misses ($12 \mid B, A \le 10000$) and Regime 2 zero-overflow sieve survival ($C_1 = 0, A \le 4000, b = 2$).}
\label{tab:beal_sieve}
\footnotesize
\begin{tabular}{@{}crrlrrrrr@{}}
\toprule
& \multicolumn{3}{c}{Regime 1: Joint CRT Near-Misses ($12 \mid B$)} & \multicolumn{5}{c}{Regime 2: Zero-Overflow Sieve ($C_1 = 0, A \le 4000$)} \\
\cmidrule(lr){2-4} \cmidrule(l){5-9}
$(x,y,z)$ & Modulus $M$ & Closest Coprime $(A, B, C)$ & Excluded by $p$ & Init $A$ & Positivity & $2$-Adic & $p \le 50$ & $p \le 181$ \\
\midrule
$(3,3,5)$ & $1,728$   & $A=7811, B=95196, C=971$       & $7, 13, 19, 31$   & $1,616$ & $152$   & $77$  & $1$  & $0$ \\
$(3,4,5)$ & $20,736$  & $A=1519, B=53148, C=6031$      & $5, 13, 19, 23$   & $1,612$ & $1,045$ & $61$  & $0$  & $0$ \\
$(4,3,5)$ & $1,728$   & $A=3463, B=105540, C=1057$     & $7, 13, 19, 31$   & $1,702$ & $24$    & $12$  & $0$  & $0$ \\
$(3,5,7)$ & $248,832$ & $A=3997, B=13331004, C=122797$ & $11, 31, 61, 101$ & $1,615$ & $1,579$ & $791$ & $15$ & $0$ \\
$(5,3,7)$ & $1,728$   & $A=1679, B=3314508, C=623$     & $13, 19, 43, 73$  & $1,613$ & $58$    & $26$  & $0$  & $0$ \\
\bottomrule
\end{tabular}
\end{table}

% ============================================================================
% BIBLIOGRAPHY
% ============================================================================
\begin{thebibliography}{11}

\bibitem{Barina2025}
D.~Barina, \emph{Improved verification limit for the convergence of the Collatz conjecture}, J. Supercomput. \textbf{81} (2025), Article 810 (verifying convergence up to $2^{71}$).

\bibitem{Bohm1978}
C.~Bohm and G.~Sontacchi, \emph{On the existence of cycles of given length in integer sequences like $x_{n+1} = x_n/2$ if $x_n$ even, and $x_{n+1} = 3x_n + 1$ otherwise}, Atti Accad. Naz. Lincei Rend. Cl. Sci. Fis. Mat. Nat. (8) \textbf{64} (1978), 260--264.

\bibitem{Eliahou1993}
S.~Eliahou, \emph{The $3x+1$ problem: new lower bounds on nontrivial cycle lengths}, Discrete Math. \textbf{118} (1993), 45--56.

\bibitem{Hercher2023}
C.~Hercher, \emph{There are no Collatz $m$-cycles with $m \le 91$}, J. Integer Seq. \textbf{26} (2023), Article 23.3.5.

\bibitem{Lagarias1985}
J.~C.~Lagarias, \emph{The $3x+1$ problem and its generalizations}, Amer. Math. Monthly \textbf{92} (1985), 3--23.

\bibitem{Laurent1995}
M.~Laurent, M.~Mignotte, and Yu.~Nesterenko, \emph{Formes lin\'{e}aires en deux logarithmes et d\'{e}terminants d'interpolation}, J. Number Theory \textbf{55} (1995), 285--321.

\bibitem{Rhin1987}
G.~Rhin, \emph{Approximants de Pad\'{e} et mesures d'irrationalit\'{e}}, Progr. Math. \textbf{71} (1987), 155--164.

\bibitem{Simons2005}
J.~Simons, \emph{On the non-existence of $2$-cycles for the $3x+1$ problem}, Math. Comp. \textbf{74} (2005), 1565--1572.

\bibitem{SimonsDeWeger2005}
J.~Simons and B.~de Weger, \emph{Theoretical and computational bounds for $m$-cycles of the $3n+1$ problem}, Acta Arith. \textbf{117} (2005), 51--70.

\bibitem{Steiner1977}
R.~P.~Steiner, \emph{A theorem on the Syracuse problem}, Proc. 7th Manitoba Conf. Numerical Mathematics and Computing (1977), 553--559.

\bibitem{Tao2022}
T.~Tao, \emph{Almost all orbits of the Collatz map attain almost bounded values}, Forum Math. Pi \textbf{10} (2022), e12, 1--56.

\bibitem{Terras1976}
R.~Terras, \emph{A stopping time problem on the positive integers}, Acta Arith. \textbf{30} (1976), 241--252.

\end{thebibliography}

\end{document}
